\chapter{南雍古典文学八讲}
\par \Large \emph{莫砺锋主编} \normalsize

\subsection*{贵妃之死：五世纪中叶南朝中国的政治、爱情、文学}

\emph{程章灿}

南朝宋孝武帝刘骏荒淫好色，宠幸殷贵妃（依谥号也称宣贵妃，一说为其堂姐妹）。殷贵妃死后文士多献诗、诔文，以谢庄《孝武帝宣贵妃诔文》最著名，但文中用“尧母之门”典，暗指殷贵妃之子刘子鸾将即位。刘骏有意传位刘子鸾，但未废当时太子刘子业。不料刘骏两年后突然死去，刘子业即位为前废帝，谢庄因此遭报复下狱，险丧命。


\subsection*{谁是唐代最伟大的诗人}

\emph{莫砺锋}

\par 《唐宋诗醇》唐代选李白、杜甫、韩愈、白居易四人。作者认为加上王维、李商隐是第一梯队。唐代最伟大诗人的观点除李杜外，清王士祯、宇文所安认为是王维\footnote{殷{\CJKfontspec{simsun.ttc}璠}《河岳英灵集》中王维、李白选诗数量相当（杜诗当时尚未成熟）。中唐时王维地位已经下降。}。日本学界长期认为是白居易（作者：理解层次不够）。林庚因特殊时代政治因素曾称是黄巢。

\par \textbf{李白与杜甫对比}：(1)对对方态度：相遇交游分别后，杜甫多次写诗怀念李白，而李白仅一首诗提到杜甫。(2)物质生活：李白尽管政治上不得志，生活上不穷困；杜甫则终生穷困。(3)思想：杜甫一生推崇儒家，而李白在儒家入世思想之外有反叛和追求自由的一面。(4)生前诗坛地位：李白生前名声很高，而杜甫并未得到同时代诗人高度评价。

\par \textbf{李杜优劣的讨论}：中唐两人都受重视，元白认为杜甫更胜一筹，韩愈认为二者同样伟大。北宋起认为杜甫地位更高，如苏轼批评李白人品上不够忠君爱国（晚年从永王军）；诗歌有时率意而锤炼不足。明高{\CJKfontspec{simsun.ttc}棅}《唐诗品汇》区分七体\footnote{五古、七古、五律、七律、五绝、七绝、五言排律。}，认为除五绝、七绝外的五体，杜甫是唯一的“大家”。新中国以来认为李杜同样伟大，如郭沫若提出李杜是“中国诗歌史上的双子星座”。作者认为杜甫地位提高的原因是对后世诗歌影响较大，而李白“靠天才作诗”，难以学习。从阅读角度，明王世贞认为“少陵稍难入，青莲较易倦”。
